%%%%%%%%%%%%%%%%%%%%%%%%%%%%%%%%%%%%%%%%%%%%%%%%%%%%%%
%
% This file defines the style for your report
% You don't need to edit it any more, if not to change the authors name.
%
% Search below for the keyword:   GROUP
% insert your group number
%
% Search below for the keyword:   AUTHORS
% insert the name of the authors
%
% If you want to compile your document you have TWO ways
% depending on the fact that 
% 	1) you have inserted only postscript images in your .tex file 
%		---> then go to MODE 1
%	2) you have inserted other kind of images (jpg, pdf, ...) in your .tex file
%		---> then go to MODE 2
%
% MODE 1 
% Type:
% 	latex homebook.tex
%
% If the compilation runs successfully and you want to see the results type:
% 	xdvi homebook.dvi &
% and use the menus to go through the document
%
% If you want to create a pdf type:
% 	dvipdfm homebook.dvi
%
% a homebook.pdf file is created
% you can see it using the command:
% 	acroread homebook.pdf &
%
%
% MODE 2
% Type:
%	pdflatex homebook.tex
%
% If the compilation runs successfully you directly have the pdf file
% and you can see it using the command:
%       acroread homebook.pdf &
%
% 
%%%%%%%%%%%%%%%%%%%%%%%%%%%%%%%%%%%%%%%%%%%%%%%%%%%%%%

\documentclass[10pt,  english, makeidx, a4paper, titlepage, oneside]{book}
\usepackage{babel}
\usepackage{fancyhdr}
\usepackage{makeidx}
\usepackage{titlesec}
\usepackage{listings}
\usepackage{booktabs}

\newenvironment{listato}{\footnotesize}{\normalsize }

%\pagestyle{empty}

\textwidth 15.5cm
\textheight 23cm
\topmargin -1cm
\oddsidemargin -0.5cm
\linespread{1.1}

\pagestyle{fancy}
\lhead{}
\chead{Microelectronic Systems}
\lfoot{}
\cfoot{}
\rfoot{}
\rhead{\thepage}

\usepackage{graphicx}
\usepackage{amsmath}
\usepackage{amsfonts}
\usepackage{amsthm}
\usepackage{amssymb}
%\oddsidemargin -1.1cm
\usepackage{graphicx}
\usepackage{caption}
\usepackage{float}
\usepackage{amsmath}
\usepackage{amssymb}
\usepackage{amsfonts}
\usepackage{amsthm}
%\usepackage{subscript}
\usepackage{empheq}
\usepackage{verbatim}
\usepackage{fancyvrb}

\lstdefinelanguage{VHDL}{morekeywords={library,use,all,entity,generic, is,port,in,out,end,architecture,of,begin,and,if,then,else,elsif,process},morecomment=[l]--}

\lstdefinestyle{vhdl}{language = VHDL, basicstyle = \ttfamily, keywordstyle = \color{keyword}\bfseries, commentstyle = \color{comment}}

\titleformat{\chapter}[display]
{\normalfont\Large\filcenter\sffamily}
{\titlerule[0.5pt]%
\vspace{1pt}
\titlerule
\vspace{1pc}
\LARGE\MakeUppercase{\chaptertitlename} \thechapter
}
{1pc}
{\titlerule
\vspace{1pc}
\Huge}

\makeindex

\begin{document}

\frontmatter
\begin{titlepage}
\vspace{0cm}
\centerline{
\includegraphics[width=3cm]{./logopoli}} 
\vspace{0.5cm}
\centerline{\LARGE Politecnico di Torino}
\vspace{2.5cm}
\centerline{\huge\sf Microelectronic Systems}
\vspace{1cm}
\centerline{\Huge\sf DLX Microprocessor: Design \& Development}
\bigskip
\centerline{\huge\sf Final Project Report}
\vspace{2cm}
\centerline{\Large Master degree in Electronics Engineering}
\bigskip
\centerline{\Large Master degree in Computer Engineering}
\vspace{4.5cm}
%%%%%%%%%%%%%%%%%%%%%%%%%%%%%%%%%%%%%%%%%%%%%%%%%%%%%%
%
\centerline{\large Referents: Prof. Mariagrazia Graziano, Giovanna Turvani}
\bigskip
\vspace{1cm}
%
%%%%%%%%%%%%%%%%%%%%%%%%%%%%%%%%%%%%%%%%%%%%%%%%%%%%%%
% GROUP
% Change the name of your group below
%
\centerline{\large Authors: group\_name}
\bigskip
%
%%%%%%%%%%%%%%%%%%%%%%%%%%%%%%%%%%%%%%%%%%%%%%%%%%%%%%
% AUTHORS
% Change the name of the Group participants here
%
\centerline{\large name1, name2}
%
%%%%%%%%%%%%%%%%%%%%%%%%%%%%%%%%%%%%%%%%%%%%%%%%%%%%%%
\vspace{2cm}
\centerline{\large \today}
\end{titlepage}

\tableofcontents

\mainmatter
%%%%%%%%%%%%%%%%%%%%%%%%%%%%%%%%%%%%%%%%%%%%%%%%%%%%%%
%    
% HERE IS WHERE YOU INCLUDE YOUR CHAPTERS
%
%%%%%%%%%%%%%%%%%%%%%%%%%%%%%%%%%%%%%%%%%%%%%%%%%%%%
% This will help you in writing your homebook
% Remember that the character % is a comment in latex
%
% chapter 1
\chapter{Prototype}
\label{chap1}
%%\graphicspath{ {./chapters/chap1images/} }
%%%%%%%%%%%%%%%%%%%%%%%%%%%%%%%%%%%%%%%%%%%%%%%%%%%%%%%%%%%
% you can organize a chapter using sections -> \section{Simulating an inverter}
% or subsections -> \subsection{simulating a particular type of inverter}
%%%%%%   First section
\section{Introduction}
In order to better and easier simulate the architecture, two scripts have been used:
A bash script to control the simulation from high level: it is in charge to invoke the C script with the right command line argument during the proper simulation phase and to 
run modelsim;
A C script in charge to perfrom various operations during different steps of the simulation, manipulating different files.
For simplicity, a single value is assigned to both operands each simulation cycle, de facto using the multiplier as a square-evaluation circuit.
\section{General simulation flow}
The C script can accepts an additional command line parameter:
-i corresponds to TEST\_INIT mode: the script reads the handwrittensamples.txt file, in which desired inputs in human-like format are stored, e.g. 12.29487, -0.9872 etc, 
and generates two other files as output, both storing data in hexadecimal encoding:
simulationinputs.hex, which contains the inputs for the Modlesim simulation
expected\_outputs.hex, which contains the outputs that are expected to be generated by the Modelsim simulation.
-v corresponds to TEST\_VALIDATION mode: the script executes a diff command between the self-generated expected\_outputs.hex file and the file generated by the Modelsim simulation using 
a pipe to execute it in background and take back its output, then processing it in order to establish if the two files are identical, Test Successful, or different, Test Failed.
The flow of the simulation is the following:
-the bash script executes the C program passing -i as command line parameter in order to generate the input vectors for the simulation and the file storing the expected results.
-Modelsim is launched and the effective output file is generated
-the C program is executed passing -v as command line parameter in order to compare the simulation results with the expected ones

\section{C script}
Since IEEE754 is the standard encoding that is normally used to store floating point data inside a flash memory, to convert a floating point value in its equivalent hexadecimal 
encoding it's sufficient to use the   \textbf{union}  C data type.
\textbf{union} data type allows to store multiple-encoding data variables in the same memory location, thus it is sufficient read a value from the handwrittensamples.hex file 
and store it as a float variable inside the previously declared memory location, and then accessing it using the other encoding
\begin{lstlisting}[style=CStyle]
    union IEEE754conv 
    {
        float f;
        uint32_t ieee754Value;
    };

\end{lstlisting}

The previously described operation is trivially implemented by the following procedure, that is very convenient both in terms synthax and in terms of CPU time, since what is done is just a couple of memory accesses:
\begin{lstlisting}[style=CStyle]
    uint32_t convFloatinIEEE754(float val)
    {
            union IEEE754conv conv;
            conv.f = val;
            return conv.ieee754Value;
    }
\end{lstlisting}

% \input{./chapters/chap_name}
% and so on
%
%%%%%%%%%%%%%%%%%%%%%%%%%%%%%%%%%%%%%%%%%%%%%%%%%%%%%%
%    
% HERE IS WHERE YOU INCLUDE YOUR APPENDICES (IF ANY)
%
\appendix
%%% Appendix A
\chapter{Adder behavioural VHDL}
\label{appendix1}

	\lstinputlisting[language=VHDL, breaklines=true]{appendices/files/adder.vhd}

% \lstinputlisting is an alternative way to import text or code from an external file. In this example the behavioural VHDL description of an adder contained in the file adder.vhd is imported. 
% Note that you can set the language of the code that you want to import (VHDL in this example). When you set the language you will see the keywords of that specific language highlighted in your output pdf file.
%You can set a lot parameters: for some examples take a look at the chapter 'How to document the project' that can you find in DLX_Project.pdf.
% \input{./appendices/appendix2}
% and so on
%
%%%%%%%%%%%%%%%%%%%%%%%%%%%%%%%%%%%%%%%%%%%%%%%%%%%%%%

\end{document}